\section{DataRobot propensity modeling and the causal weighting contract}
\label{sec:exec25-datarobot-reuse}
DataRobot's readmission accelerator provides a concrete reference for separating
prediction of treatment from evaluation of an intervention. The retained
official page describes propensity modeling followed by inverse probability
of treatment weighting. The linked notebook was acquired at revision
\texttt{fd2670c8b18380e712ba}\allowbreak\texttt{64ed3fe3ea055c114b4e},
dated 26 March 2025. This is the notebook-path revision returned by the bounded
7 October 2026 cutoff query, not a claim about the latest entire repository.
Its bytes match the original Git tree; code and explanatory cells were
inspected without running code or treating saved outputs as reproduced results.
The official overview's later acquisition date and displayed September 2026
update do not authenticate its exact cutoff bytes.~\citep{exec13_ch_datarobot,exec25_datarobot_notebook}

For the banking proposal, the corresponding question might concern whether
an account receives a servicing call and whether it becomes delinquent within
30 days. A predictive model of call assignment estimates a propensity under
the observed assignment regime; it is not itself an intervention effect.
Identification requires a treatment definition and population, consistency,
appropriate exchangeability given eligible pretreatment information, positivity,
and an observation law supporting the intended uncertainty calculation.
Features measured after the call or caused by it cannot be admitted merely
because they improve treatment prediction. In the inspected notebook, removing
the readmission outcome leaves the other variables available to the propensity
model. That code operation does not establish their causal timing for a new
banking question. The proposal therefore retains a separate reviewed covariate
and clock contract.~\citep{exec25_datarobot_notebook,exec25_datarobot_choice}

\subsection{What the inspected notebook actually does}
The selected workflow creates a DataRobot project, models the medication-status
target, waits for the modeling process and retrieves a top model. It inspects
predictive metrics including AUC and log loss, then requests predictions on the
feature rows from the study. Original explanatory text explicitly calls these
in-sample predictions and mentions a separate training-predictions API as an
alternative. That statement is preserved as the author's description; it is
not adopted as a general justification of inference with fitted propensity
scores. The correct use of sample splitting, cross-fitting or another
finite-sample/asymptotic argument depends on the estimator and sampling law.
Neither a high AUC nor low log loss alone establishes covariate balance,
positivity, causal identification or calibrated uncertainty.~\citep{exec25_datarobot_notebook,exec25_datarobot_choice}

The notebook selects the predicted \texttt{class\_Yes} column, indexes it by
\texttt{row\_id}, and joins it to input rows numbered from zero. A proposed
adapter must verify class meaning, unique row identity, ordering, completeness
and finite probability range before constructing weights. A missing row after
a left join must not disappear from the denominator. The artifact contract
must identify the project, model, prediction job, feature schema and exact
input rows; a generic prediction table is insufficient to recover the
assignment probability for a different treatment or data version. These are
proposed conditions for reuse, not claims of a demonstrated upstream mismatch.
\citep{exec25_datarobot_notebook,exec25_datarobot_choice}

The selected weight function computes, for binary treatment $T_i$ and fitted
propensity $\widehat e_i$, the expression
\begin{equation}
 w_i=\frac{T_i\widehat p}{\widehat e_i}
       +\frac{(1-T_i)(1-\widehat p)}{1-\widehat e_i},
\end{equation}
where $\widehat p$ is the empirical treated fraction. Its non-stabilized option
replaces the two numerators by one. The inspected function has no clipping or
zero-denominator branch. This is a static scope statement, not an observed
numerical failure. A production design needs an explicit invalid-input and
positivity rule. Clipping or trimming can be useful design choices, but they
must have their own target, bias and uncertainty analysis rather than being
added as invisible numerical repairs.~\citep{exec25_datarobot_notebook,exec25_datarobot_choice}

\subsection{Stabilization does not remove a normalized arm's extreme weights}
The notebook's categorical comparison sums weights by treatment and category,
then normalizes within the treatment arm. Multiplying every weight in one arm
by the same positive factor cancels from this normalized percentage. For an
exact authored treated-arm example, let propensities be $1/4$ and $3/4$, and
outcomes be one and zero. Unstabilized weights are $4$ and $4/3$. With treated
fraction $1/2$, stabilized weights are $2$ and $2/3$. Both normalized means are
\begin{equation}
 \frac{4\cdot1+(4/3)\cdot0}{4+4/3}
 =\frac{2\cdot1+(2/3)\cdot0}{2+2/3}=\frac34.
\end{equation}
The commonly used weight-concentration diagnostic
$\operatorname{ESS}=(\sum_i w_i)^2/\sum_i w_i^2$ is likewise unchanged at
$8/5$ in this two-row example. Stabilization changes absolute weights; it does
not improve their relative concentration within a normalized arm by itself.
This calculation concerns that displayed normalized functional, not every
possible weighting estimator or uncertainty procedure.~\citep{exec25_datarobot_notebook,exec25_datarobot_choice}

With the same stabilization numerator and a propensity of $1/100$, a treated
row receives weight 50. Replacing that propensity by $1/10$ would reduce the
weight to five, but this is an explicitly authored hypothetical clipping
example, not notebook behavior or a justified scientific correction. In a
servicing study, such an account may represent a rarely treated stratum whose
counterfactual outcome is weakly supported by the data. A useful product
should expose that weakness, possible target restriction and sensitivity,
rather than hiding it behind a stable-looking plot. A weighted outcome contrast
still requires an uncertainty method accounting for fitted weights and the
actual independent, clustered, panel or longitudinal sampling structure.
\citep{exec25_datarobot_choice}

The retained notebook's final comparison displays normalized weighted
readmission percentages. It discusses both a possible treatment effect and
unobserved confounding, and invites additional assumptions and information.
The review does not attribute the saved numerical conclusion to a new run.
A chart of two percentages is not, by itself, a predeclared effect estimator
with calibrated confidence limits. The proposed banking contract must name
the effect scale, target population, selection rule, nuisance-fitting strategy,
variance or resampling law and sensitivity analysis. In-sample prediction is
not declared universally invalid, but any resulting inferential guarantee must
be justified for the selected design.~\citep{exec25_datarobot_notebook,exec25_datarobot_choice}

\subsection{Reference choice, rights and public qualification}
The selected disposition is reference use of the documented workflow and
inspected notebook. The retained complete revision tree contains no root
license, and the pinned license API returned 404. Public availability does
not establish permission to copy, port or deploy the accelerator. The
commercial DataRobot environment, API entitlement, data rights and dependency
versions are separate unresolved items. The notebook reads a local modified
clinical dataset; this review neither acquires nor executes that dataset, and
the scholarly citations inside the notebook are not newly admitted as locally
inspected papers. Direct adoption and source porting remain deferred pending
rights and an exact supported environment. Conditional integration could use
an entitled propensity service with an independently justified causal
estimator and conversion contract. Rejection is limited to treating predictive
quality or weighted plots as sufficient proof of causal validity. It does not
claim that other DataRobot offerings lack causal functionality.
\citep{exec13_ch_datarobot,exec25_datarobot_notebook,exec25_datarobot_choice}

A proposed public study would compare a fixed conventional propensity model,
a permitted frozen DataRobot model and a proposed contract-wrapped workflow on
the same generated assignment/outcome laws and held-out tasks. Known
propensities and effects provide exact small-model oracles; permitted public
observational data supplies workflow stress without invented causal truth.
Freeze feature timing, target, class mapping, task split, model selection,
weight rules, overlap criteria, uncertainty law and error tolerances before
measurement. Fixtures should include an exchanged class label, duplicate or
missing prediction rows, probabilities at the boundary, near-positivity
violations, post-treatment covariates, failed prediction jobs and an output
from the wrong model or question version. All failures and scientifically
valid refusals remain in the request denominator. No such native comparison
has been conducted in this reference decision.~\citep{exec25_datarobot_choice}

Scientific owners would approve the question, covariates, weights and
uncertainty; adapter owners would verify class and row alignment; platform and
rights owners would select an entitled model/API environment and permitted
data; evaluators would control generated truth and held-out scoring;
maintainers would review model and version drift. These are proposed roles.
Hypothetical preparation of USD 1,000, all model attempts of USD 2,400,
scientific review of USD 2,800 and maintenance allocation of USD 1,000 total
USD 7,200. At 120 attempts, 90 completions and 60 admissible outputs, costs are
USD 60, USD 80 and USD 120 respectively. These exact authored denominators do
not establish price or investment value. The conditional route requires
permitted artifacts, justified causal inference and matched complete costs.
\citep{exec25_datarobot_choice}

\subsection{From a component choice to useful assistance}

The notebook comparison illustrates the chapter's acquisition rule. Predictive
quality can make a propensity component useful, while the causal analysis
still needs eligible information, a defined population, support and justified
uncertainty. Stabilizing weights does not itself change their concentration
within a normalized arm. A successful modeling job does not settle source
permission, conversion correctness or the cost of obtaining an admissible
result. Those are distinct reasons to prefer a bounded bridge, an existing
workflow, another component or no integration.

The same standard applies when language assistance is added. An assistant can
help formulate a request, explain assumptions and operate reviewed tools, but
its contribution must be compared on the complete scientific task and cost.
The next part examines that contribution, retaining unresolved information
and adverse outcomes as part of the proposed service.
